% !TEX root = ../main.tex
% ---------------------------------------------------------------------
%  Docker -- Software-Virtualisierung mit Containern
% ---------------------------------------------------------------------
\section{Docker}
\label{sec:docker}

\SW{Wozu} Plattform zur \T{Software-Virtualisierung}: Software wird in
standardisierte, wiederverwendbare Einheiten verpackt, die \emph{alles}
enthalten, was zum Ausführen nötig ist. Diese Einheiten laufen auf nahezu
jedem Rechner direkt \ra portabel, skalierbar, unkompliziert.
Das Problem, das sie löst: „Bei mir läuft's!“ -- Abhängigkeiten,
Versionskonflikte, Installationsprozeduren.

\SW{Abgrenzung} Docker virtualisiert \emph{nur den Linux-Kernel}, nicht
die ganze Maschine. Deshalb gilt immer \T{Container-OS $=$ Host-OS} --
anders als bei einer virtuellen Maschine. Auf Windows/macOS braucht Docker
deshalb eine VM (WSL). Die Programme laufen direkt auf dem Host-Kernel;
in \R{ps ax} auf dem Host sieht man den Container-Prozess.

\subsection{Image vs.\ Container}
\begin{ztab}{@{}lXX@{}}
 & Image & Container \\ \midrule
Was & Bauplan, alles Nötige zum Ausführen & ein \emph{ausgeführtes} Image \\
Dateisystem & \T{read-only}, unveränderlich & \T{read-write} Overlay über dem Image \\
Anzahl & eines & beliebig viele aus einem Image \\
Platz & volle Grösse & nur die \emph{Änderungen} gegenüber dem Image \\
Lebensdauer & bis \R{rmi} & bis \R{rm}, überlebt das Stoppen \\
\end{ztab}

\begin{formel}[Der Satz, der alles trägt]
Das Ausführen eines \emph{Images} erzeugt einen \emph{Container} --
nie umgekehrt. Änderungen zur Laufzeit (Logfiles, DB-Einträge,
nachinstallierte Pakete) landen im Overlay des Containers, das Image
bleibt unberührt.
\end{formel}

\subsection{Host und Kapselung}
\SW{Host} der Rechner mit dem OS, auf dem die Container laufen (hier
Raspberry Pi OS). Er stellt nur die grundlegendsten OS-Funktionen bereit:
nacktes Dateisystem, Prozesse, Sockets. \emph{Nicht} Software, Daten,
User/Gruppen oder IP-Ports -- das muss der Container selbst mitbringen.

\SW{Kapselung} Aus dem Container heraus normalerweise kein Zugriff auf das
Dateisystem des Hosts oder anderer Container; gilt auch für User und Ports.
Der Host dagegen kann Container und Images ändern und löschen.
Die zwei Löcher in dieser Wand sind \R{docker cp} und \T{Volumes}.

\subsection{Lebenszyklus}
Vier Zustände: \T{created} $\to$ \T{running} $\to$ \T{stopped} $\to$
\T{deleted}, dazu das lokale Image-Repository als fünfter Knoten:
\begin{lstlisting}
Docker Hub --pull--> lokales Image-Repo
  lokales Repo --run-----> running
  lokales Repo --create--> created --start--> running
  running  --stop--> stopped --start/restart--> running
  stopped/created --rm--> deleted
  Container --commit--> lokales Image-Repo
\end{lstlisting}
\begin{ztab}{@{}lX@{}}
\R{docker run <img>} & Image \ra Container erzeugen \emph{und} starten \\
\R{docker create <img>} & nur erzeugen (Zustand \emph{created}) \\
\R{docker start <cnt>} & starten, \emph{detached} (Hintergrund) \\
\R{docker start -i <cnt>} & starten, \emph{attached} (Vordergrund) \\
\R{docker stop <cnt>} & running \ra stopped \\
\R{docker restart <cnt>} & stop, dann start \\
\R{docker rm <cnt>} & Container löschen \\
\R{docker rmi <img>} & Image löschen \\
\end{ztab}

\begin{falle}[Drei Stolperstellen]
1.~\T{Jedes \R{run} erzeugt einen \emph{neuen} Container} -- nicht
\R{run} zum Wiederstarten nehmen, sondern \R{start}.\\
2.~Ein Image lässt sich \textbf{nicht} löschen, solange noch ein Container
davon existiert.\\
3.~Wer \R{-it} beim \R{create} vergisst, kann den Container später nicht
interaktiv nutzen -- die Option muss schon beim \emph{Erzeugen} dabei sein.
\end{falle}

\subsection{Die wichtigsten Kommandos}
\begin{ztab}{@{}lX@{}}
\R{docker image ls} & lokale Images (kurz: \R{docker images}) \\
\R{docker ps} & laufende Container (lang: \R{container ls}) \\
\R{docker ps -a} & auch beendete \\
\R{docker ps -a --size} & mit Platzbedarf \\
\R{docker run -it <img>} & interaktiv (siehe unten) \\
\R{docker run --rm <img>} & Container nach Ende sofort löschen \\
\R{docker run --name <n> <img>} & eigener Name statt Zufallsname \\
\R{docker rename <alt> <neu>} & nachträglich umbenennen \\
\R{docker attach <cnt>} & Terminal an Hintergrund-Container hängen \\
\R{docker container prune} & \emph{alle} gestoppten Container löschen \\
\R{docker login -u <user>} & an Docker Hub anmelden \\
\R{docker image push <tag>} & Image auf Docker Hub laden \\
\end{ztab}
Docker vergibt jedem Container eine zufällige \T{ID} und einen zufälligen
\T{Namen}; Kommandos akzeptieren beides. Fehlt ein Image lokal, lädt
\R{run} es automatisch von Docker Hub.

\begin{formel}[Was \R{-i} und \R{-t} bedeuten]
\R{-i} verbindet den Container mit der Standardeingabe, \R{-t} legt einen
Pseudo-Terminalemulator davor. Zusammen \R{-it} ergibt eine
Terminal-Sitzung wie bei \R{ssh}.\\
\T{Ohne \R{-it} beendet sich ein Shell-Container sofort}: die \R{bash}
findet keine Standardeingabe und bricht ab -- der häufigste Anfängerfehler.
\end{formel}

\subsection{Beenden und Trennen}
\begin{ztab}{@{}lX@{}}
\R{exit} & beendet die Shell im Container \ra Container \emph{stopped} \\
\R{<Ctrl-d>} & dasselbe (beendet die Shell) \\
\R{<Ctrl-c>} & \emph{nicht-interaktiv}: beendet \R{docker run} und damit
  den Container; \emph{interaktiv}: beendet nur das gerade laufende
  Programm \emph{in} der Shell, nicht die Shell \\
\R{<Ctrl-p> <Ctrl-q>} & \T{detach}: nur das Terminal trennen, Container
  läuft weiter. Nur bei Containern mit \R{-it}. Zurück mit \R{attach}. \\
\end{ztab}
In allen Fällen \emph{existiert} der Container weiter (Status
\emph{Exited}), bis man ihn mit \R{rm} löscht.

\subsection{Schlanke Images}
Images sind bewusst minimal -- im Ubuntu-Image fehlen \R{ping}, \R{nano},
\R{wget}, \R{curl}; vorhanden sind \R{cp}, \R{mv}, \R{mkdir}, \R{cat},
\R{tar}, \R{apt}, \R{ps}, \R{top}.
\SW{Nachinstallieren} erst die Paketlisten aktualisieren, dann installieren:
\begin{lstlisting}
apt -y update && apt -y install nano
\end{lstlisting}
\R{-y} beantwortet alle Rückfragen mit Ja (nötig, weil beim Bauen von
Images niemand tippt). Die Verkettung mit \R{\&\&} ist üblich -- warum,
zeigt sich beim Dockerfile.

\SW{Auch die Prozesse sind minimal} Auf einem normalen Linux laufen viele
Dienste (\R{sshd}, \R{dockerd}, \R{dhcpd}, \dots{} -- in Unix
\emph{daemons}, Endung \R{d}). Im Ubuntu-Container läuft \T{genau ein
Prozess: \R{bash}} -- der Container soll nur eine Software ausführen, den
Rest stellt der Host bereit. Zählen mit \R{ps -aux | wc -l}.

\begin{bsp}[Platzbedarf -- typische Prüfungsfrage]
\R{docker ps -a --size} zeigt \R{0B (virtual 192MB)}: der Container selbst
belegt nichts, solange er nichts geändert hat; die 192\,MB sind das Image.
Installiert man darin 20\,MB Software, wächst der \emph{Container} um
20\,MB, das Image bleibt gleich gross.\\
\ra 4 Container aus einem 50-MB-Image ohne Änderungen brauchen zusammen
\T{50\,MB}, nicht 200. Und: Dateien \emph{im Container} zu löschen gibt
keinen Platz frei, wenn sie aus dem Image stammen.
\end{bsp}

\subsection{Hintergrund und zusätzliche Prozesse}
\begin{ztab}{@{}lX@{}}
\R{docker run -d <img>} & \T{detached} starten; läuft, bis der gestartete
  Prozess sich beendet. Für Serverdienste ohne Shell-Interaktion. \\
\R{docker run -itd <img>} & interaktiv \emph{und} detached; kehrt sofort
  zurück, später \R{attach} \\
\R{docker logs <cnt>} & Ausgabe des Containers ansehen \\
\R{docker logs -f <cnt>} & \emph{folgt} der Ausgabe wie \R{tail -f};
  Abbruch mit \R{<Ctrl-c>} \\
\R{docker exec <cnt> <cmd>} & \T{zusätzlichen} Prozess im laufenden
  Container starten \\
\R{docker exec -it <cnt> bash} & zweite Shell parallel -- das
  Debug-Werkzeug \\
\R{docker exec -d <cnt> <cmd>} & dieser Prozess im Hintergrund \\
\R{docker top <cnt> ax} & Prozesse \emph{im} Container auflisten \\
\end{ztab}

\begin{formel}[\R{run} gegen \R{exec} -- der Unterschied]
\R{run} erzeugt einen \T{neuen Container} aus einem Image.
\R{exec} startet einen \T{weiteren Prozess} in einem \emph{bereits
laufenden} Container. \R{exec} auf einen gestoppten Container
funktioniert nicht.
\end{formel}

\subsection{Dateien kopieren: \R{docker cp}}
\begin{lstlisting}
docker cp <cnt>:<quelle> <ziel>    # Container -> Host
docker cp <quelle> <cnt>:<ziel>    # Host -> Container
\end{lstlisting}
Funktioniert auch bei einem Container im Zustand \emph{created}, also
bevor er je lief. Typische Abfolge: \R{docker create --name myubuntu -it
ubuntu} $\to$ \R{docker cp hello.txt myubuntu:/tmp/} $\to$
\R{docker start myubuntu} $\to$ \R{docker attach myubuntu}.

\subsection{Das geschichtete Container-Dateisystem}
Das virtuelle Dateisystem eines Containers besteht aus übereinander
gelegten Schichten:
\begin{ztab}{@{}lX@{}}
oben & \T{Container-Layer}, read-write -- die Änderungen \\
Mitte & weitere \emph{Image}-Schichten: nachinstallierte Software,
  Webserver, DB, dotnet \dots{} \\
unten & \T{Base-Image}, meist ein minimales Linux (Ubuntu, Debian) \\
\end{ztab}
Alle Schichten ausser der obersten sind \T{read-only}.
\T{Jeder Container bekommt seinen eigenen Container-Layer} -- auch mehrere
aus demselben Image. Der Layer wird \emph{mit} dem Container gelöscht.
\merkzeile Praktisch für reinen Tisch (temporäre Dateien verschwinden),
unpraktisch für Software, die Daten dauerhaft speichern muss.

\SW{Wo das physisch liegt} \R{/var/lib/docker/overlay2/\dots} und
\R{/var/lib/containerd/\dots}. Suchen mit
\R{sudo find /var/lib/ -name <datei>}.

\begin{falle}[Zwei Beobachtungen dazu]
1.~Bei einem \emph{laufenden} Container findet \R{find} die Datei
\T{zweimal} (docker und containerd), bei einem gestoppten nur einmal.\\
2.~Diese Verzeichnisse werden von Docker verwaltet und dürfen
\textbf{nie} direkt verändert werden.
\end{falle}

\subsection{Volumes}
Ein Volume bildet einen Pfad des Host-Dateisystems in das
Container-Dateisystem ab. Änderungen sind auf beiden Seiten sofort
sichtbar, und die Dateien \T{überleben das Löschen des Containers} --
alles andere im Container-Layer nicht.

\begin{ztab}{@{}lX@{}}
\T{Host Volume} & \R{-v /pfad/im/host:/pfad/im/cnt} -- Nutzer bestimmt
  beide Seiten \\
\T{Anonymes V.} & \R{-v /pfad/im/cnt} -- Docker legt das Host-Verzeichnis
  unter \R{/var/lib/docker/volumes/<hash>/\_data/} an; kaum gebraucht \\
\T{Named Volume} & \R{-v name:/pfad/im/cnt} -- wie anonym, aber mit Namen \\
\end{ztab}

\begin{bsp}[Wofür Volumes wirklich da sind]
Versions-Update eines Datenbankservers: Die containerisierte DB legt ihre
Daten im Volume ab. Beim Update wird der alte Container gestoppt und
gelöscht -- die Daten bleiben. Der neue Container aus dem aktuellen Image
greift auf dieselben Daten zu.
\end{bsp}

\SW{Administration} \R{docker volume ls} $\cdot$
\R{docker volume rm <vol>} $\cdot$ \R{docker volume prune}.
Volumes überleben ihren Container, müssen also von Hand aufgeräumt werden.

\begin{falle}[Warum \R{volume prune} trügerisch ist]
Es löscht \emph{alle} Volumes, die gerade keinem Container zugeordnet
sind. Genau die wertvollen Daten eines Volumes, dessen Container du
versehentlich schon gelöscht hast, sind damit weg -- und es gibt
\textbf{kein} Docker-Kommando, um vorher den Inhalt eines Volumes
anzusehen.
\end{falle}

\subsection{Image aus Container: \R{docker commit}}
\begin{lstlisting}
docker commit <container> <image>
\end{lstlisting}
Friert die Änderungen eines laufenden oder gestoppten Containers als neue
Image-Schicht ein. Damit lässt sich manuell installierte Software in ein
Image übernehmen: Ubuntu starten $\to$ \R{apt install nano} $\to$
\R{commit} $\to$ das neue Image enthält nano.
\merkzeile Das ist die Antwort auf „Kann aus einem Container ein Image
erzeugt werden?“ -- ja, mit \R{commit}. Der saubere Weg bleibt aber das
Dockerfile.

\subsection{Repositories und Docker Hub}
\T{Docker Hub} ist das Standard-Repository: \R{docker run} und
\R{docker pull} holen Images von dort. Öffentliche Repositories sind für
jeden zugänglich, private nur für den Besitzer (bei Docker Hub ab einer
gewissen Anzahl kostenpflichtig) -- oder on premise mit eigener
Repository-Software (CNCF Distribution).

\begin{ztab}{@{}lX@{}}
\R{docker search <begriff>} & Images suchen; Spalten \R{STARS},
  \R{OFFICIAL}, \R{AUTOMATED}. Offizielle Images sind vertrauenswürdiger. \\
\R{docker pull <img>} & herunterladen, ohne zu starten \\
\end{ztab}

\begin{formel}[Aufbau eines Image-Namens]
\begin{center}\ttfamily
<repo>/<source>/<image>:<tag>
\end{center}
\begin{tabularx}{\linewidth}{@{}lX@{}}
\R{<image>} & der Name, z.\,B. \R{ubuntu}, \R{hello-world} \\
\R{<source>} & Nutzer oder Organisation auf Docker Hub;
  Standard \R{library} (die offiziellen Images) \\
\R{<tag>} & die Version; Standard \R{latest} \\
\R{<repo>} & die Registry; Standard \R{index.docker.io} \\
\end{tabularx}\\[2pt]
Also ist \R{docker run hello-world} identisch mit\\
\R{docker run index.docker.io/library/hello-world:latest}.\\
Weitere Beispiele: \R{ubuntu:jammy} (22.04),
\R{kaohslu/epp-demo}, \R{mcr.microsoft.com/dotnet/runtime:5.0}.
\end{formel}

\begin{falle}[\R{run} aktualisiert nicht]
\R{docker run <img>} nimmt die \T{neueste lokal vorhandene} Version --
auch wenn im Repository längst eine neuere liegt. Nur \R{docker pull}
holt die aktuelle.\\
Und beim \R{pull}: Das alte Image wird \emph{nicht} überschrieben,
sondern verliert nur seinen Tag und erscheint mit \R{<none>}
(\emph{dangling image}) -- Container könnten es noch brauchen.
\end{falle}

\subsection{Der Weg eines eigenen Images}
\begin{lstlisting}
docker build -t <user>/<name> .   # Image bauen und taggen
docker login -u <user>            # an Docker Hub anmelden
docker image push <user>/<name>   # hochladen
docker run <user>/<name>          # auf beliebigem Host starten
\end{lstlisting}
Drei Zeilen zum Veröffentlichen, eine zum Benutzen -- ohne dass auf dem
Zielrechner irgendeine Laufzeitumgebung installiert sein muss.
